\section{Metodologia}\label{sec:metodologia}

\subsection{Erros de seguimento}

Com o modelo cinemático de bicicleta ($\dot\psi = v\tan\delta/L$), os dois controladores usam erros definidos no centro do eixo dianteiro $p_f$ \cite{hoffmann2007stanley}. Sendo $P$ o ponto da referência mais próximo de $p_f$, $\psi_r$ a direção da sua tangente e $\hat n_r$ a normal à esquerda,
\begin{equation}
e_y = (P - p_f)\cdot \hat n_r, \qquad e_\psi = \psi_r - \psi .
\end{equation}
Assim, $e_y>0$ indica referência à esquerda e $e_\psi>0$ indica que o veículo precisa girar para a esquerda.

\subsection{Controlador on-off}

Um on-off que atuasse apenas sobre $e_\psi$ deixaria o veículo paralelo à referência, sem eliminar o deslocamento lateral. Por isso, a referência de orientação foi definida como a direção de um ponto da referência a uma distância $L_d=0{,}8\,\mathrm{m}$ adiante de $P$, ao longo da tangente. No referencial do veículo, esse ponto está exatamente no ângulo
\begin{equation}
e = e_\psi + \arctan\left(e_y/L_d\right),
\label{eq:erro-onoff}
\end{equation}
que é o erro usado pelo on-off. O esterçamento assume três níveis, com histerese:
\begin{equation}
\delta_k = +\delta_{on}\ \text{se}\ e>h_{on};\quad -\delta_{on}\ \text{se}\ e<-h_{on};\quad 0\ \text{se}\ |e|<h_{off};\quad \delta_{k-1}\ \text{nos demais casos}.
\end{equation}
O nível zero evita oscilação contínua entre os batentes nas retas. A amplitude $\delta_{on}$ deve superar o esterçamento de regime $\arctan(L\kappa)$ da curva mais fechada, mas, quanto maior, maior o ciclo-limite; o valor de $15^\circ$ foi escolhido pelos ensaios de sensibilidade (\autoref{sec:resultados}). O comando passa pelo mesmo limitador de taxa do Stanley.

\subsection{Controlador Stanley}

O Stanley \cite{thrun2006stanley,hoffmann2007stanley,snider2009steering} soma a correção da orientação a um termo não linear do erro lateral:
\begin{equation}
\delta = \operatorname{sat}_{\pm\delta_{\max}}\left[e_\psi + \arctan\left(\frac{k\,e_y}{v_s+|v|}\right)\right].
\label{eq:stanley}
\end{equation}
Para erros pequenos e sem saturação, $\dot e_y \approx -k|v|\,e_y/(v_s+|v|)$: convergência exponencial com constante de tempo próxima de $1/k$, aproximadamente independente da velocidade. Com $k=1{,}5\,\mathrm{s^{-1}}$ e $v=1\,\mathrm{m/s}$, ela vale $\approx0{,}87\,\mathrm{s}$.

\subsection{Referências geométricas}

Na primeira etapa, os erros foram calculados com a pose verdadeira do simulador em relação a: uma \textbf{reta} de $28\,\mathrm{m}$ deslocada $0{,}5\,\mathrm{m}$ para a esquerda da pose inicial (convergência a partir de um erro inicial); um \textbf{círculo} de raio $3\,\mathrm{m}$ tangente à pose inicial (curvatura constante); e um \textbf{oito} formado por dois círculos de $3\,\mathrm{m}$ percorridos em sentidos opostos, em cujo cruzamento a curvatura salta de $+1/3$ para $-1/3\,\mathrm{m^{-1}}$. Os ensaios foram feitos numa região plana e sem cones, para a qual o carrinho é reposicionado (a cena não é salva), com $v_{ref}=1{,}0$ e $2{,}0\,\mathrm{m/s}$.

\subsection{Estimação da referência por visão computacional}\label{sec:visao}

Na segunda etapa, $e_y$ e $e_\psi$ são estimados somente a partir da imagem, a cada passo:
\begin{enumerate}
    \item \textbf{Segmentação}: conversão para HSV e seleção de $40^\circ\le H\le75^\circ$, $S\ge0{,}4$, $V\ge0{,}2$, o que elimina o asfalto e as faixas brancas.
    \item \textbf{Centro da linha}: a máscara é percorrida de baixo para cima; em cada linha da imagem é escolhido o segmento amarelo mais próximo do centro anterior, descartando saltos maiores que $20\%$ da largura (outros trechos da linha ao fundo).
    \item \textbf{Projeção no chão}: cada centro gera um raio pelo modelo \textit{pinhole}, intersectado com o plano da pista.
    \item \textbf{Memória curta}: pontos dos últimos $3\,\mathrm{s}$ a até $3\,\mathrm{m}$ do eixo dianteiro são mantidos, para cobrir a região cega; ela só é usada se houve detecção nos últimos $0{,}3\,\mathrm{s}$.
    \item \textbf{Ajuste local}: no referencial do veículo, o trecho até $\approx1{,}2\,\mathrm{m}$ à frente é ajustado por uma parábola, com rejeição de pontos atípicos; o ponto mais próximo do eixo dianteiro fornece $e_y$, $e_\psi$ e a curvatura $\kappa$. Em curvas fechadas, um arco de circunferência substitui a parábola quando tem menor resíduo, pois extrapola melhor até o eixo.
\end{enumerate}
A medida é rejeitada com menos de seis pontos, extensão menor que $0{,}3\,\mathrm{m}$ ou resíduo acima de $12\,\mathrm{cm}$. Numa falha curta, o último esterçamento é mantido por até $0{,}4\,\mathrm{s}$; depois, a velocidade de referência é zerada. A projeção e a memória usam a pose da câmera dada pelo simulador (equivalente a odometria ideal), mas os erros enviados ao controle vêm apenas dos pixels detectados.

\subsection{Ensaios na pista e critérios de avaliação}

Na pista, $v_{ref}$ foi mantida constante ($0{,}5$ a $2{,}5\,\mathrm{m/s}$, o máximo do veículo) para isolar o efeito da velocidade. Também foi testada a redução nas curvas, $v_{ref}^{*}=\min\{v_{ref},\max[0{,}5;\ \min(\sqrt{a_{lat}/|\kappa|},\ v_{ref}/(1+2|e_y|+|e_\psi|))]\}$, com $a_{lat}=1{,}5\,\mathrm{m/s^2}$ e $\kappa$ estimada pela câmera. Cada ensaio parte da pose original da cena e termina ao completar a volta (retorno a menos de $1\,\mathrm{m}$ do início) ou numa falha: inclinação da carroceria acima de $45^\circ$ (capotamento), distância maior que $1\,\mathrm{m}$ da linha (saída da pista) ou mais de $1\,\mathrm{s}$ sem detecção (perdido).

Como o erro estimado pela câmera não é independente, foi calculado também um erro verdadeiro: a distância do eixo dianteiro ao esqueleto da linha amarela extraído da textura da pista (pixels de $1{,}8\,\mathrm{cm}$), usado apenas na avaliação e nos critérios de término. As métricas desconsideram os primeiros $5\,\mathrm{s}$.
